\section{The base system: GeFL-F}\label{ch:gefl}

Model-heterogeneous methods differ in what the clients exchange: predictions on shared inputs (FedDF, FedMD), class
means of internal representations (FedProto, FedTGP), a shared part of the network (LG-FedAvg, FedGH, FedMRL), or a
generator of synthetic data (FedGen, GeFL). The paper builds on GeFL-F, the feature-level variant of GeFL. It splits
every client's network into two parts, a small \emph{extractor} that all clients share and a \emph{head} that is
specific to the client's architecture, and lets the clients share a generator of extractor outputs. Training runs in
three stages (\cref{fig:s-pipeline}).

\begin{figure}[t]
\centering
\begin{tikzpicture}[font=\small,
  box/.style={draw,rounded corners=2pt,align=center,minimum height=1.05cm,text width=3.9cm,fill=gray!6},
  arr/.style={-{Stealth[length=2mm]},thick}]
\node[box] (s1) {\textbf{Stage I} ($T_{\mathrm I}$ rounds)\\ extractor $\phi$ + heads, FedAvg;\\ then freeze $\phi$};
\node[box,right=0.6cm of s1] (s2) {\textbf{Stage II} ($T_{\mathrm{II}}$ rounds)\\ generator $G(h\mid y)$, FedAvg\\ (Adam, coupled decay)};
\node[box,right=0.6cm of s2] (s3) {\textbf{Stage III} ($T_{\mathrm{III}}$ rounds)\\ each head: $T_s$ epochs generated,\\ then $T_r$ epochs real};
\draw[arr] (s1) -- node[above]{$h=\phi(x)$} (s2);
\draw[arr] (s2) -- node[above]{$G$} (s3);
\node[below=0.35cm of s2,align=center,text width=12.5cm,font=\footnotesize] {Every client sends parameters to the
server in each round of every stage. Stage I averages the extractor over all clients (and each head over the clients of
its architecture); Stage II averages the generator over all clients; Stage III averages each head over the clients of
its architecture, which does nothing when every architecture has one client.};
\end{tikzpicture}
\caption{GeFL-F's three stages. The extractor and the generator are shared by all clients; the heads are not.}
\label{fig:s-pipeline}
\end{figure}

\paragraph{Stage I: the shared extractor} All clients train a small convolutional extractor $\phi$ together with
their heads by FedAvg, for $T_{\mathrm I}=50$ rounds on CIFAR. The extractor is the same network for everyone, so it
can be averaged. On CIFAR it is two $3\times3$ convolutions with three and ten output channels, each followed by
batch normalisation and ReLU, and a $2\times2$ max pooling. A $32\times32$ image becomes a $10\times16\times16$ array,
the \emph{feature} $h=\phi(x)$, with $d=2560$ numbers. After Stage I the extractor is frozen, and every client computes
the features of its images once, without augmentation. From then on, every client works on features, not images.

\paragraph{The heads} Client $k$ trains a head of its own architecture, one of ten convolutional networks CNN-1 to
CNN-10 with one to four convolutional blocks of different widths. Each head has two parts. Its \emph{body}, the
convolutional blocks, maps a feature $h$ to an \emph{embedding} $e$; its \emph{last layer} is linear and maps $e$ to
the $C$ logits $f=We+b$. In the main setting there are ten clients and client $k$ trains architecture $k$, so each
architecture has exactly one client.

\paragraph{Stage II: the shared generator} The clients train a conditional variational autoencoder $G(h\mid y)$ on
their features, again by FedAvg, for $T_{\mathrm{II}}=200$ rounds, with the Adam optimiser and coupled weight decay
(an L2 penalty $\lambda\|\theta\|^2/2$ added to the loss). The generator takes a class label $y$ and random noise $z$
and outputs a feature $\tilde h$ that should look like a real feature of class $y$. The label enters through a linear
projection of the one-hot vector $e_y$; the weights that multiply class $c$ form a vector $r_c$, the \emph{class row}.
The class rows are the only parameters of the generator that belong to one class: everything the generator knows that
distinguishes class $c$ from the others must pass through $r_c$.

\paragraph{Stage III: training the heads} For $T_{\mathrm{III}}=100$ rounds, every client trains its head in two
phases. In the \emph{synthetic phase} it draws labels uniformly over all $C$ classes, generates features for them, and
trains for $T_s$ epochs on these pairs with cross-entropy ($T_s=1$ in GeFL-F). In the \emph{real phase} it trains for
$T_r=5$ epochs on its own features. The server then averages the heads of each architecture. With one client per
architecture this averaging does nothing, so in the main setting the only knowledge a head receives from other
clients comes through the generator. This fact returns in \cref{ch:results,ch:defence}.

\paragraph{Why a generator helps} Client 2 of the running example holds no image of class 9. Its real phase can only
teach its head to separate the classes it holds. The synthetic phase can show the head features of class 9, made by a
generator that was trained partly on client 7's 63 images of that class. If the generated class-9 features look like
real ones, client 2's head learns class 9 without ever seeing a class-9 image. Whether they do look like real ones is
the central question of \cref{ch:diagnosis}.

\paragraph{What goes wrong} GeFL-F was designed and evaluated on balanced clients. Under the long tail of the running
example it reaches \R{main/cifar10/GeFL-F/bal} balanced accuracy on CIFAR-10-LT, about the level of Local, in which
every client trains alone (\R{main/cifar10/Local/bal}). On the tail classes, the 40\% least frequent (classes 6--9),
its recall is \R{main/cifar10/GeFL-F/tail}. On CIFAR-100-LT it reaches \R{main/cifar100/GeFL-F/bal}, against
\R{main/cifar100/Local/bal} for Local. Federation through the generator brings almost nothing.

Where, inside these three stages, are the rare classes lost?
